% !TEX program = xelatex
% 完整中文理论稿；数学结构、标签与公式编号对应 main.tex。
\documentclass[11pt,a4paper]{article}
\usepackage[margin=1in]{geometry}
\usepackage{amsmath,amsthm,mathtools}
\usepackage{fontspec}
% Native XeTeX line breaking avoids a locally incompatible xeCJK release.
\XeTeXlinebreaklocale "zh"
\XeTeXlinebreakskip=0pt plus 1pt
\setmainfont{SimSun}[BoldFont=SimHei,ItalicFont=SimSun,BoldItalicFont=SimHei]
\newfontfamily\LatinText{Times New Roman}
\newcommand{\LatinWord}[1]{{\LatinText #1}}
\usepackage{unicode-math}
\setmathfont{Cambria Math}
\usepackage[expansion=false]{microtype}
\usepackage{algorithm,algpseudocode}
\renewcommand{\abstractname}{摘要}
\renewcommand{\refname}{参考文献}
\renewcommand{\proofname}{证明}
\floatname{algorithm}{算法}
\algrenewcommand\algorithmicrequire{\textbf{输入：}}
\algrenewcommand\algorithmicensure{\textbf{输出：}}
\algrenewcommand\algorithmicreturn{\textbf{返回}}
\algrenewcommand\algorithmicwhile{\textbf{当}}
\algrenewcommand\algorithmicdo{\textbf{时}}
\algrenewcommand\algorithmicif{\textbf{如果}}
\algrenewcommand\algorithmicthen{\textbf{则}}
\algrenewtext{EndWhile}{\textbf{结束循环}}
\algrenewtext{EndIf}{\textbf{结束判断}}
\usepackage{flafter,placeins}
\usepackage[hidelinks]{hyperref}
\theoremstyle{plain}
\newtheorem{proposition}{命题}[section]
\pdfstringdefDisableCommands{\def\LatinWord#1{#1}}
\hypersetup{pdftitle={Wasserstein 鲁棒假设检验中的变量选择},
 pdfauthor={研究笔记},
 pdfsubject={有限检验问题、仿射对偶表示与切平面}}
\setlength{\parindent}{1.3em}
\setlength{\parskip}{0pt}
\setlength{\emergencystretch}{2em}
\clubpenalty=10000
\widowpenalty=10000
\setlength{\jot}{2pt}
\allowdisplaybreaks[1]
\AtBeginDocument{%
\setlength{\abovedisplayskip}{6pt plus 2pt minus 1pt}%
\setlength{\belowdisplayskip}{6pt plus 2pt minus 1pt}}
\title{\LatinWord{Wasserstein} 鲁棒假设检验中的变量选择}
\author{}
\date{2026年10月18日}
\makeatletter
\renewcommand{\maketitle}{%
\begin{center}
{\Large\bfseries\@title\par}\vspace{4pt}
{\small\@date\par}
\end{center}\vspace{2pt}}
\makeatother

\begin{document}
\maketitle
\begin{abstract}
Wasserstein 鲁棒检验的变量选择问题具有精确有限表示及最大仿射对偶表示。
对偶耦合生成全局割。理论分析给出由求解精度和坐标影响界决定的误差界与迭代次数界。
\end{abstract}

\section{问题设定}\label{sec:problem}
令 $[D]=\{1,\ldots,D\}$，$0\le d\le D$ 为整数。
对于可测评分函数 $f:\mathbb R^D\to\mathbb R$，定义
\begin{equation}\label{eq:risk}
R(f)=\sup_{\substack{P_0\in\mathcal U_0\\P_1\in\mathcal U_1}}
\left\{\mathbb E_{P_0}\ell(-f(X))+
       \mathbb E_{P_1}\ell(f(X))\right\}.
\end{equation}
损失 $\ell$ 有限、非负、凸且非减，并满足 $\ell(0)>0$。
记 $\ell_0(s)=\ell(-s)$，$\ell_1(s)=\ell(s)$。
对于 $S\subseteq[D]$，以 $\operatorname{Proj}_S:\mathbb R^D\to\mathbb R^{|S|}$
表示坐标投影，并定义
\[
\mathcal G_S=\{g:\mathbb R^{|S|}\to\mathbb R:\ g\text{ 为Borel可测}\}.
\]
变量选择问题为
\begin{equation}\label{eq:problem}
\boxed{V(S)=\inf_{g\in\mathcal G_S}R(g\circ\operatorname{Proj}_S),
\qquad \mathrm{OPT}=\min_{|S|\le d}V(S).}
\end{equation}

$N$ 个观测为 $x_1,\ldots,x_N\in\mathbb R^D$。
非空索引集 $I_0,I_1$ 按类别划分 $[N]$，类内概率质量为 $q_i>0$：
\[
\sum_{i\in I_k}q_i=1,\qquad
\widehat P_k=\sum_{i\in I_k}q_i\delta_{x_i},\qquad
\mathcal U_k=\{P:W_1(P,\widehat P_k)\le\theta_k\},\quad k=0,1.
\]
其中 $\delta_x$ 为点 $x$ 处的单位质量，$\theta_k>0$，
$W_1$ 对具有有限一阶矩的分布使用 $\ell_1$ 代价。
不另加权时，$i\in I_k$ 的样本取 $q_i=1/|I_k|$。

\section{对偶重构}\label{sec:value}
固定 $S$，记
\[
x_i^S=\operatorname{Proj}_Sx_i,\qquad \Delta_{ir,j}=|x_{ij}-x_{rj}|,\qquad
c_{ir}^S=\|x_i^S-x_r^S\|_1=\sum_{j\in S}\Delta_{ir,j}.
\]
其中 $i,r\in[N]$，$j\in[D]$。
耦合 $\pi_{ir}\ge0$ 将质量从源样本 $i$ 运输到目标样本 $r$。定义
\[
\mathcal T_k(S)=\left\{\pi\ge0:
\sum_r\pi_{ir}=q_i,\quad
\sum_{i\in I_k,r}\pi_{ir}c_{ir}^S\le\theta_k\right\}.
\]
以 $u=(u_1,\ldots,u_N)$ 为目标点评分变量，有限运输表示为
\begin{equation}\label{eq:finite-transport}
V(S)=\inf_{u\in\mathbb R^N}
\sum_{k=0}^1\max_{\pi\in\mathcal T_k(S)}
\sum_{i\in I_k,r}\pi_{ir}\ell_k(u_r).
\end{equation}
对各运输线性规划取对偶，引入源样本损失上界 $a_i\in\mathbb R$
和运输价格 $\lambda_k\ge0$。
记 $J(a,\lambda)=\sum_iq_i a_i+\sum_k\theta_k\lambda_k$，得到
\begin{equation}\label{eq:finite}
\begin{aligned}
V(S)=\min_{a,u\in\mathbb R^N,\ \lambda\ge0}\quad&J(a,\lambda)\\
\text{要求}\quad&
\ell_k(u_r)\le a_i+\lambda_kc_{ir}^S,
\quad i\in I_k,\ r\in[N],\ k=0,1.
\end{aligned}
\end{equation}
附录~\ref{app:finite} 证明精确性，并恢复最优检验。

对于二元指示向量 $z=\mathbf1_S$，用坐标价格 $w_{kj}$ 代替按选择开关的价格 $\lambda_kz_j$。
由于 $a_i\ge0$，
\[
\sum_k\theta_kw_k=z\sum_k\theta_k\lambda_k\le zJ(a,\lambda)
\quad\text{当 }w_k=\lambda_kz.
\]
固定 $z\in[0,1]^D$，定义
\begin{equation}\label{eq:value}
\begin{aligned}
F(z)=\min_{a,u,\lambda,w}\quad&J(a,\lambda)\\
\text{要求}\quad&
\ell_k(u_r)\le a_i+\sum_j\Delta_{ir,j}w_{kj},
\quad i\in I_k,\ r\in[N],\ k=0,1,\\
&0\le w_k\le\lambda_k\mathbf1,\qquad k=0,1,\\
&\theta_0w_0+\theta_1w_1\le zJ(a,\lambda),\\
&\lambda\ge0,\qquad a,u\in\mathbb R^N.
\end{aligned}
\end{equation}
其中 $w_k=(w_{k1},\ldots,w_{kD})^\top$，$\mathbf1=(1,\ldots,1)^\top$；
向量不等式按坐标理解。对于二元 $z$，在其余约束下，
\[
\theta_0w_0+\theta_1w_1\le zJ
\quad\Longleftrightarrow\quad
w_{kj}=0\quad\text{当 }z_j=0,\ k=0,1.
\]
当 $z_j=1$ 时，该坐标上的这条不等式冗余，可以省略：
\[
\theta_0w_{0j}+\theta_1w_{1j}
\le\sum_k\theta_k\lambda_k\le J.
\]
删除整组约束则也会移除未选坐标上的限制，使每个 $z$ 都对应 $V([D])$。
保留它则有 $F(\mathbf1_S)=V(S)$；命题~\ref{prop:dual}\textup{(i)} 证明投影后的可行域相同，且目标保持不变。

式~\eqref{eq:value} 的对偶使用经缩放的运输耦合 $\Pi\ge0$、
标量 $\alpha\ge0$ 及坐标乘子 $\mu\ge0$，其形式为
\begin{equation}\label{eq:dual-normalized}
\begin{aligned}
F(z)=\max_{\Pi,\alpha,\mu}\quad&H(\Pi)\\
\text{要求}\quad&(\Pi,\alpha,\mu)\in\mathcal D,\qquad
\alpha+\mu^\top z=1.
\end{aligned}
\end{equation}
其中
\[
H(\Pi)=\sum_{r=1}^N\inf_{s\in\mathbb R}
\sum_{k=0}^1\sum_{i\in I_k}\Pi_{ir}\ell_k(s),
\]
并以 $(v)_+=\max\{v,0\}$ 记正部，有
\begin{equation}\label{eq:dual-domain}
\mathcal D=\left\{(\Pi,\alpha,\mu)\ge0:
\begin{array}{l}
\displaystyle\sum_r\Pi_{ir}=\alpha q_i,\quad i\in[N],\\[3pt]
\displaystyle\sum_j
\left(\sum_{i\in I_k,r}\Pi_{ir}\Delta_{ir,j}-\theta_k\mu_j\right)_+
\le\alpha\theta_k,\quad k=0,1
\end{array}\right\}.
\end{equation}
该耦合的源质量为 $\alpha q_i$；$H$ 为每个目标点选取一个评分，
使其加权损失最小。拉格朗日推导见命题~\ref{prop:dual}\textup{(i)}。

对于归一化三元组 $\alpha+\mu^\top z=1$，
将三个对偶变量同时乘以 $\tau\ge0$ 仍属于 $\mathcal D$，并有
\[
H(\tau\Pi)=\tau H(\Pi),\qquad
\max_{\tau\ge0}\{2\sqrt{\tau H(\Pi)}-\tau\}=H(\Pi),
\quad \tau^*=H(\Pi).
\]
在每条射线上优化尺度，可消去归一化约束：
\begin{equation}\label{eq:dual-affine}
\boxed{F(z)=\max_{(\Pi,\alpha,\mu)\in\mathcal D}
\{2\sqrt{H(\Pi)}-\alpha-\mu^\top z\}.}
\end{equation}
命题~\ref{prop:dual}\textup{(ii)} 证明这一等价性及最优解的存在性。归一化最优解 $(\bar\Pi,\bar\alpha,\bar\mu)$ 对应的仿射最优解为 $F(z)(\bar\Pi,\bar\alpha,\bar\mu)$；割的斜率也必须乘以这一尺度。

\section{切平面算法}\label{sec:algorithm}
令 $\mathcal Z_d=\{z\in\{0,1\}^D:\mathbf1^\top z\le d\}$。
在 $z^t$ 处用式~\eqref{eq:dual-affine} 评估 $F(z^t)$，
取最优解 $(\Pi^t,\alpha^t,\mu^t)$，得到在 $z^t$ 处取等的全局仿射下界：
\begin{equation}\label{eq:tangent}
\begin{aligned}
c_t(z)&:=2\sqrt{H(\Pi^t)}-\alpha^t-(\mu^t)^\top z,\\
F(z)&\ge c_t(z)=F(z^t)-(\mu^t)^\top(z-z^t),\qquad
c_t(z^t)=F(z^t).
\end{aligned}
\end{equation}
下一选集和下界由以下问题给出
\begin{equation}\label{eq:master}
\begin{aligned}
(z^{t+1},\eta^{t+1})\in
\operatorname*{arg\,min}_{z\in\mathcal Z_d,\ \eta\ge0}\quad&\eta\\
\text{要求}\quad&
\eta\ge F(z^s)-(\mu^s)^\top(z-z^s),
\qquad s=1,\ldots,t.
\end{aligned}
\end{equation}
为给出迭代次数界，定义

\begin{equation}\label{eq:slope-bound}
C=2\ell(0),\qquad
A_{kj}=\sum_{i\in I_k}q_i\max_{r\in[N]}\Delta_{ir,j},\qquad
M_j=C\max_{k=0,1}\frac{A_{kj}}{\theta_k}.
\end{equation}

将这些界排序为 $M_{(1)}\ge\cdots\ge M_{(D)}$，对 $m=0,\ldots,D$ 令

\begin{equation}\label{eq:rate-quantities}
\rho_m=\sum_{r=m+1}^{\min(D,m+d)}M_{(r)},\qquad
K_m=\sum_{r=0}^{\min(d,m)}\binom mr.
\end{equation}

其中 $\rho_m$ 控制前 $m$ 个坐标之外至多 $d$ 个坐标的总影响，
$K_m$ 是前 $m$ 个坐标上可能选法的数量。

对非精确求解，保留全部全局仿射下界 $c_s(z)\le F(z)$，以 $p_s$ 记对应的原始可行值。定义

\[
\begin{gathered}
U_t=\min_{s\le t}p_s,\qquad
m_t(z)=\max\{0,c_1(z),\ldots,c_t(z)\},\\
\underline L_t\le\min_{z\in\mathcal Z_d}m_t(z),\qquad
L_t=\max\{0,\underline L_1,\ldots,\underline L_t\},\qquad
G_t=U_t-L_t.
\end{gathered}
\]

因此 $\underline L_t$ 是主问题的全局下界，$G_t$ 是计算得到的最优性间隙。
以 $f_t$ 记 $U_t$ 对应的恢复评分函数。

\begin{proposition}[收敛性]\label{prop:convergence}
\textup{(i)} 对每个 $z\in[0,1]^D$，在仿射对偶中加入 $\alpha\le C$、
$0\le\mu_j\le M_j$ 后，其最优值仍为 $F(z)$ 且可以取到。

\textup{(ii)} 固定 $\delta,\beta\ge0$。设返回的仿射下界斜率为 $-\mu^s$，且每次内层求解和主问题求解均满足

\[
\begin{gathered}
0\le\mu^s\le M,\qquad 0\le p_s-c_s(z^s)\le\delta,\\
z^{t+1}\in\mathcal Z_d,\qquad
\eta^{t+1}\ge m_t(z^{t+1}),\qquad
0\le\eta^{t+1}-L_t\le\beta.
\end{gathered}
\]

对每个已完成的迭代 $t$ 及满足 $t\ge K_m$ 的 $m$，有

\begin{equation}\label{eq:certificate}
0\le R(f_t)-\mathrm{OPT}\le G_t\le\delta+\beta+\rho_m.
\end{equation}

当 $\varepsilon\ge\delta+\beta$ 且 $\varepsilon>0$ 时，
以 $G_t\le\varepsilon$ 为准则停止所需的迭代次数满足

\begin{equation}\label{eq:iteration-bound}
T_\varepsilon\le K_{m_\varepsilon},\qquad
m_\varepsilon=\min\{m\in\{0,\ldots,D\}:\rho_m\le\varepsilon-\delta-\beta\}.
\end{equation}

精确求解对应 $\delta=\beta=0$、$L_t=\eta^{t+1}$。
证明见附录~\ref{app:convergence}。
\end{proposition}

\begin{algorithm}[H]
\caption{变量选择的精确切平面算法}\label{alg:oa}
\begin{algorithmic}[1]
\Require 问题数据、预算 $d$、容差 $\varepsilon>0$ 和 $z^1\in\mathcal Z_d$。
\Ensure 选集 $S^*$、其最优检验 $f^*$，以及宽度至多为 $\varepsilon$ 的风险区间。
\State $t\gets1$，$U\gets+\infty$。
\While{真}
  \State 在 $z^t$ 处精确求解式~\eqref{eq:dual-affine}；按命题~\ref{prop:convergence}\textup{(i)} 选取满足 $0\le\mu_j^t\le M_j$ 的 $\mu^t$。
  \If{$F(z^t)<U$}
    \State $(U,z^*)\gets(F(z^t),z^t)$。
  \EndIf
  \State 保留索引为 $t$ 的割~\eqref{eq:tangent}。
  \State 精确求解式~\eqref{eq:master}，得到 $(z^{t+1},\eta^{t+1})$。
  \If{$U-\eta^{t+1}\le\varepsilon$}
    \State $S^*\gets\{j:z_j^*=1\}$。
    \State 对 $S^*$ 求解式~\eqref{eq:finite}，按式~\eqref{eq:extension} 构造 $g^*$，并令 $f^*=g^*\circ\operatorname{Proj}_{S^*}$。
    \State \Return $(S^*,f^*,[\eta^{t+1},U])$。
  \EndIf
  \State $t\gets t+1$。
\EndWhile
\end{algorithmic}
\end{algorithm}

\FloatBarrier
\appendix
\section{有限表示的证明}\label{app:finite}
设 $x_i\in\mathbb R^D$，$i\in[N]$，非空集合 $I_0,I_1$ 划分 $[N]$。
设 $q_i>0$，$\sum_{i\in I_k}q_i=1$，$\theta_k>0$。
损失 $\ell:\mathbb R\to[0,\infty)$ 有限、凸且非减，满足 $\ell(0)>0$；
记 $\ell_0(s)=\ell(-s)$，$\ell_1(s)=\ell(s)$。
对于 $S\subseteq[D]$，以 $\operatorname{Proj}_S$ 表示坐标投影，
以 $\delta_x$ 表示点 $x$ 处的单位质量。定义
\[
\begin{gathered}
x_i^S=\operatorname{Proj}_Sx_i,\qquad
c_{ir}^S=\|x_i^S-x_r^S\|_1,\qquad
\widehat P_k=\sum_{i\in I_k}q_i\delta_{x_i},\\
\mathcal U_k=\{P:W_1(P,\widehat P_k)\le\theta_k\},\qquad
\mathcal G_S=\{g:\mathbb R^{|S|}\to\mathbb R:\ g\text{ 为Borel可测}\},\\
R(f)=\sum_{k=0}^1\sup_{P\in\mathcal U_k}\mathbb E_P\ell_k(f(X)),
\qquad
V(S)=\inf_{g\in\mathcal G_S}R(g\circ\operatorname{Proj}_S).
\end{gathered}
\]
其中 $W_1$ 对有限一阶矩分布使用 $\ell_1$ 运输代价。设
\begin{equation}\label{eq:appendix-finite}
\begin{aligned}
T(S)=\inf_{a,u\in\mathbb R^N,\ \lambda\in\mathbb R_+^2}\quad&
J(a,\lambda):=\sum_iq_i a_i+\sum_{k=0}^1\theta_k\lambda_k\\
\text{要求}\quad&
\ell_k(u_r)\le a_i+\lambda_kc_{ir}^S,
\quad i\in I_k,\ r\in[N],\ k=0,1.
\end{aligned}
\end{equation}
在证明与 $V(S)$ 相等之前，以 $T(S)$ 单独表示有限问题的值。

\begin{proposition}\label{prop:finite}
对于每个 $S\subseteq[D]$，有 $V(S)=T(S)$，定义 $T(S)$ 的下确界可以取到。
每个有限最优解都可构造一个 Borel 评分函数
$g^*:\mathbb R^{|S|}\to\mathbb R$，满足
$R(g^*\circ\operatorname{Proj}_S)=V(S)$。
\end{proposition}
\begin{proof}
\textit{有限运输对偶。}
固定 $u\in\mathbb R^N$，定义
\[
\mathcal T_k(S)=\left\{\pi\in\mathbb R_+^{|I_k|\times N}:
\sum_r\pi_{ir}=q_i\ (i\in I_k),\quad
\sum_{i\in I_k,r}\pi_{ir}c_{ir}^S\le\theta_k\right\}.
\]
取 $a_i\in\mathbb R$、$\lambda_k\ge0$，以下 $i\in I_k$、$r\in[N]$：
\begin{equation}\label{eq:transport-lagrangian}
\begin{aligned}
L_k(\pi,a,\lambda_k;u)
={}&\sum_{i,r}\pi_{ir}\ell_k(u_r)
+\sum_i a_i\left(q_i-\sum_r\pi_{ir}\right)
+\lambda_k\left(\theta_k-\sum_{i,r}\pi_{ir}c_{ir}^S\right)\\
={}&\sum_iq_i a_i+\theta_k\lambda_k
+\sum_{i,r}\pi_{ir}[\ell_k(u_r)-a_i-\lambda_kc_{ir}^S].
\end{aligned}
\end{equation}
\[
\sup_{\pi\ge0}L_k=
\begin{cases}
\sum_iq_i a_i+\theta_k\lambda_k,
&\ell_k(u_r)\le a_i+\lambda_kc_{ir}^S\quad\forall i,r,\\
+\infty,&\text{其他情形}.
\end{cases}
\]
由线性规划强对偶，
\begin{equation}\label{eq:transport-dual}
\begin{aligned}
\max_{\pi\in\mathcal T_k(S)}\sum_{i\in I_k,r}\pi_{ir}\ell_k(u_r)
=\min_{(a_i)_{i\in I_k},\ \lambda_k\ge0}\quad&
\sum_{i\in I_k}q_i a_i+\theta_k\lambda_k\\
\text{要求}\quad&
\ell_k(u_r)\le a_i+\lambda_kc_{ir}^S,\quad i\in I_k,\ r\in[N].
\end{aligned}
\end{equation}
\[
T(S)=\inf_{u\in\mathbb R^N}\sum_{k=0}^1
\max_{\pi\in\mathcal T_k(S)}\sum_{i\in I_k,r}\pi_{ir}\ell_k(u_r).
\]
这里只对有限线性规划取对偶，未交换极小与极大的顺序。

\textit{下界：$T(S)\le V(S)$。}
对 $g\in\mathcal G_S$、$\pi\in\mathcal T_k(S)$，取
$\widetilde x_{ir}\in\mathbb R^D$ 满足
\[
\operatorname{Proj}_S\widetilde x_{ir}=x_r^S,\qquad
\operatorname{Proj}_{S^c}\widetilde x_{ir}=\operatorname{Proj}_{S^c}x_i,
\qquad \|\widetilde x_{ir}-x_i\|_1=c_{ir}^S.
\]
\begin{equation}\label{eq:expectation-transport}
\begin{aligned}
P_k^\pi&=\sum_{i\in I_k,r}\pi_{ir}\delta_{\widetilde x_{ir}},
&W_1(P_k^\pi,\widehat P_k)&\le\sum_{i\in I_k,r}\pi_{ir}c_{ir}^S\le\theta_k,\\
\mathbb E_{P_k^\pi}\ell_k(g(\operatorname{Proj}_S X))
&=\sum_{i\in I_k,r}\pi_{ir}\ell_k(g(x_r^S)).
\end{aligned}
\end{equation}
\begin{equation}\label{eq:finite-lower}
\begin{aligned}
V(S)&\ge\inf_{g\in\mathcal G_S}\sum_k
\max_{\pi\in\mathcal T_k(S)}\sum_{i\in I_k,r}\pi_{ir}\ell_k(g(x_r^S))\\
&\ge\inf_{u\in\mathbb R^N}\sum_k
\max_{\pi\in\mathcal T_k(S)}\sum_{i\in I_k,r}\pi_{ir}\ell_k(u_r)=T(S).
\end{aligned}
\end{equation}
\textit{上界：$V(S)\le T(S)$。}
对任意有限可行解 $(a,u,\lambda)$ 与 $y\in\mathbb R^{|S|}$，令
\[
\begin{gathered}
E_k(y)=\min_{i\in I_k}\{a_i+\lambda_k\|x_i^S-y\|_1\},\qquad
i_k(y)=\min\operatorname*{arg\,min}_{i\in I_k}
\{a_i+\lambda_k\|x_i^S-y\|_1\},\\
r_k=\|x_{i_k}^S-y\|_1,\qquad
\rho=\|x_{i_0}^S-x_{i_1}^S\|_1\le r_0+r_1.
\end{gathered}
\]
\begin{equation}\label{eq:extension}
g(y)=\begin{cases}
(1-\gamma)u_{i_0}+\gamma u_{i_1},\quad
\gamma=\min\{1,r_0/\rho\},&\rho>0,\\
u_{i_0},&\rho=0.
\end{cases}
\end{equation}
当 $\rho>0$ 时，$\gamma\rho\le r_0$、$(1-\gamma)\rho\le r_1$。
由凸性与有限可行性，
\[
\begin{aligned}
\ell_0(g(y))&\le(1-\gamma)\ell_0(u_{i_0})+\gamma\ell_0(u_{i_1})
\le a_{i_0}+\lambda_0\gamma\rho\le E_0(y),\\
\ell_1(g(y))&\le(1-\gamma)\ell_1(u_{i_0})+\gamma\ell_1(u_{i_1})
\le a_{i_1}+\lambda_1(1-\gamma)\rho\le E_1(y).
\end{aligned}
\]
当 $\rho=0$ 时，源与目标对 $(i_0,i_0)$、$(i_1,i_0)$ 给出相同的两个界。
最小指标规则保证 $g$ 为 Borel 可测，且 $|g(y)|\le\max_r|u_r|$；无须插值所有 $u_r$。
对 $P\in\mathcal U_k$，将最优耦合分解为源测度 $\nu_i$（重复源按 $q_i$ 拆分）：
\[
\nu_i(\mathbb R^D)=q_i,\qquad
\sum_{i\in I_k}\nu_i=P,\qquad
\sum_{i\in I_k}\int\|x-x_i\|_1\,\nu_i(dx)\le\theta_k.
\]
\begin{equation}\label{eq:extension-risk}
\begin{aligned}
\mathbb E_P\ell_k(g(\operatorname{Proj}_S X))
&\le\sum_{i\in I_k}\int
[a_i+\lambda_k\|\operatorname{Proj}_Sx-x_i^S\|_1]\,\nu_i(dx)\\
&\le\sum_{i\in I_k}q_i a_i+\theta_k\lambda_k.
\end{aligned}
\end{equation}
\[
R(g\circ\operatorname{Proj}_S)\le J(a,\lambda),\qquad V(S)\le T(S).
\]
可行点 $(a,u,\lambda)=(\ell(0)\mathbf1,0,0)$ 的目标为 $C=2\ell(0)$。
在闭子水平集 $J\le C$ 上，
\[
0\le a_i\le C/q_i,\qquad 0\le\lambda_k\le C/\theta_k.
\]
损失约束界定 $\ell(u_r)$ 与 $\ell(-u_r)$。
若 $\ell$ 非常数，则 $t\to\infty$ 时 $\ell(t)\to\infty$，故 $u$ 有界；
若 $\ell$ 为常数，则限制 $u=0$。
由 Weierstrass 定理存在最优解，其延拓达到 $V(S)$。
\end{proof}

\section{对偶重构的证明}\label{app:dual}
设 $x_i\in\mathbb R^D$，非空集合 $I_0,I_1$ 划分 $[N]$，且
$q_i,\theta_k>0$、$\sum_{i\in I_k}q_i=1$。
设 $\ell:\mathbb R\to[0,\infty)$ 有限、凸且单调不减，
满足 $\ell(0)>0$；记 $\ell_0(s)=\ell(-s)$、$\ell_1(s)=\ell(s)$。
令
\[
\Delta_{ir,j}=|x_{ij}-x_{rj}|,\qquad
c_{ir}^S=\sum_{j\in S}\Delta_{ir,j},\qquad
J(a,\lambda)=\sum_iq_i a_i+\sum_{k=0}^1\theta_k\lambda_k.
\]
对 $S\subseteq[D]$ 与 $z\in[0,1]^D$，考虑
\begin{equation}\label{eq:appendix-fixed}
\begin{aligned}
V(S)=\min_{a,u,\lambda}\quad&J(a,\lambda)\\
\text{要求}\quad&
\ell_k(u_r)\le a_i+\lambda_kc_{ir}^S,
\quad i\in I_k,\ r\in[N],\ k=0,1,\\
&a,u\in\mathbb R^N,\qquad\lambda\in\mathbb R_+^2,
\end{aligned}
\end{equation}
\begin{equation}\label{eq:appendix-extended}
\begin{aligned}
F(z)=\min_{a,u,\lambda,w}\quad&J(a,\lambda)\\
\text{要求}\quad&
\ell_k(u_r)\le a_i+\sum_{j=1}^D\Delta_{ir,j}w_{kj},
\quad i\in I_k,\ r\in[N],\ k=0,1,\\
&0\le w_{kj}\le\lambda_k,
\quad k=0,1,\ j\in[D],\\
&\sum_{k=0}^1\theta_kw_{kj}\le z_jJ(a,\lambda),
\quad j\in[D],\\
&a,u\in\mathbb R^N,\qquad\lambda\in\mathbb R_+^2.
\end{aligned}
\end{equation}
对 $\Pi\in\mathbb R_+^{N\times N}$，定义
\[
H(\Pi)=\sum_{r=1}^N\inf_{s\in\mathbb R}
\sum_{k=0}^1\sum_{i\in I_k}\Pi_{ir}\ell_k(s),\qquad
d_{kj}(\Pi)=\sum_{i\in I_k,r}\Pi_{ir}\Delta_{ir,j}.
\]
\begin{equation}\label{eq:appendix-cone}
\mathcal D=\left\{(\Pi,\alpha,\mu)\ge0:
\begin{array}{l}
\sum_r\Pi_{ir}=\alpha q_i\quad(i\in[N]),\\[2pt]
\sum_j(d_{kj}(\Pi)-\theta_k\mu_j)_+\le\alpha\theta_k
\quad(k=0,1)
\end{array}\right\},\qquad(v)_+=\max\{v,0\}.
\end{equation}

\begin{proposition}[对偶表示与次梯度]\label{prop:dual}
\textup{(i)} 对每个 $S\subseteq[D]$，有 $F(\mathbf1_S)=V(S)$。
对每个 $z\in[0,1]^D$，原问题的最小值与归一化对偶的最大值均取得，且
\begin{equation}\label{eq:appendix-normalized}
0<F(z)=\max_{\substack{(\Pi,\alpha,\mu)\in\mathcal D\\
\alpha+\mu^\top z=1}}H(\Pi)\le2\ell(0).
\end{equation}

\textup{(ii)} 相同的值也由以下仿射表示取得：
\begin{equation}\label{eq:appendix-affine}
F(z)=\max_{(\Pi,\alpha,\mu)\in\mathcal D}
\{2\sqrt{H(\Pi)}-\alpha-\mu^\top z\}.
\end{equation}
对任意在 $z$ 处的归一化最优解 $(\bar\Pi,\bar\alpha,\bar\mu)$，
以下三元组是仿射最优解：
\begin{equation}\label{eq:dual-rescaling}
(\Pi,\alpha,\mu)=F(z)(\bar\Pi,\bar\alpha,\bar\mu),\qquad
H(\Pi)=F(z)^2.
\end{equation}

\textup{(iii)} 函数 $F$ 是凸的。对任意在 $z^t$ 处的仿射最优解
$(\Pi^t,\alpha_t,\mu^t)$，其仿射函数满足
\begin{equation}\label{eq:support-contact}
\begin{aligned}
c_t(z)&:=2\sqrt{H(\Pi^t)}-\alpha_t-(\mu^t)^\top z
\le F(z)\quad(z\in[0,1]^D),\\
c_t(z^t)&=F(z^t),\qquad
c_t(z)=F(z^t)-(\mu^t)^\top(z-z^t).
\end{aligned}
\end{equation}
因此 $-\mu^t\in\partial F(z^t)$，其中次梯度按
$[0,1]^D$ 上的不等式理解；不要求可微。
\end{proposition}
\begin{proof}
\textup{(i)} 对角约束给出 $a_i\ge\ell_k(u_i)\ge0$。
对式~\eqref{eq:appendix-fixed} 的可行三元组，令
$z=\mathbf1_S$、$w_k=\lambda_kz$，则
\[
\ell_k(u_r)\le a_i+\sum_j\Delta_{ir,j}w_{kj},\qquad
\sum_k\theta_kw_{kj}=z_j\sum_k\theta_k\lambda_k\le z_jJ(a,\lambda).
\]
反之，式~\eqref{eq:appendix-extended} 在 $z=\mathbf1_S$ 处给出
\[
\begin{gathered}
w_{kj}=0\quad(j\notin S),\qquad0\le w_{kj}\le\lambda_k\quad(j\in S),\\
\ell_k(u_r)\le a_i+\sum_{j\in S}\Delta_{ir,j}w_{kj}
\le a_i+\lambda_kc_{ir}^S.
\end{gathered}
\]
两个映射均保持 $J$，故二元情形等价。

对式~\eqref{eq:appendix-extended} 的损失、价格上限及选择约束，
分别配置乘子 $\Pi_{ir},\beta_{kj},\mu_j\ge0$。
保留 $a,u$ 自由、$\lambda,w\ge0$；令 $\alpha=1-\mu^\top z$。
拉格朗日函数为
\begin{equation}\label{eq:lagrangian}
\begin{aligned}
\mathcal L={}&J+\sum_{k,i\in I_k,r}\Pi_{ir}
\left[\ell_k(u_r)-a_i-\sum_j\Delta_{ir,j}w_{kj}\right]
+\sum_{k,j}\beta_{kj}(w_{kj}-\lambda_k)\\
&+\sum_j\mu_j\left[\sum_k\theta_kw_{kj}-z_jJ\right]\\
={}&\sum_i a_i\left(\alpha q_i-\sum_r\Pi_{ir}\right)
+\sum_{r,k,i\in I_k}\Pi_{ir}\ell_k(u_r)\\
&+\sum_k\lambda_k\left(\alpha\theta_k-\sum_j\beta_{kj}\right)
+\sum_{k,j}w_{kj}[\beta_{kj}+\theta_k\mu_j-d_{kj}(\Pi)].
\end{aligned}
\end{equation}
对 $a,\lambda,w$ 的下确界有限当且仅当
\begin{equation}\label{eq:dual-conditions}
\sum_r\Pi_{ir}=\alpha q_i,\qquad
\sum_j\beta_{kj}\le\alpha\theta_k,\qquad
\beta_{kj}\ge(d_{kj}(\Pi)-\theta_k\mu_j)_+.
\end{equation}
行和约束强制 $\alpha\ge0$。按目标列分别取下确界，得
\begin{equation}\label{eq:dual-loss-inf}
\begin{aligned}
\inf_{u\in\mathbb R^N}
\sum_{r=1}^N\sum_{k=0}^1\sum_{i\in I_k}\Pi_{ir}\ell_k(u_r)
&=\sum_{r=1}^N\inf_{s\in\mathbb R}
\sum_{k=0}^1\sum_{i\in I_k}\Pi_{ir}\ell_k(s)\\
&=H(\Pi).
\end{aligned}
\end{equation}
消去 $\beta$ 是精确的，因为其最小可行分量为
$\beta_{kj}=(d_{kj}(\Pi)-\theta_k\mu_j)_+$。
由此得到 $\mathcal D$ 与 $\alpha+\mu^\top z=1$。

对 $J_0=\{j:z_j=0\}$，可行性强制 $w_{kj}=0$。
消去这些变量及选择约束后，一个严格可行点为
\[
u=0,\qquad a_i=\ell(0)+1,\qquad\lambda_k=1,\qquad
w_{kj}=z_j/2\quad(j\notin J_0).
\]
损失约束与价格上限约束均严格成立，且
\[
\sum_k\theta_kw_{kj}=\frac{z_j}{2}\sum_k\theta_k
<z_j\left[2(\ell(0)+1)+\sum_k\theta_k\right]=z_jJ
\quad(j\notin J_0).
\]
由 Slater 定理，约化问题强对偶且对偶最大值取得，
这也包含 $J_0=[D]$ 的情形。对消去的坐标补回
\[
\beta_{kj}=0,\qquad
\mu_j=\max_{k=0,1}d_{kj}(\Pi)/\theta_k\quad(j\in J_0).
\]
这些取值满足式~\eqref{eq:dual-conditions}，且既不改变目标，
也不改变归一化等式。

常数可行点 $(a,u,\lambda,w)=(\ell(0)\mathbf1,0,0,0)$
的目标为 $C=2\ell(0)$。其闭下水平集 $J\le C$ 满足
\[
\begin{gathered}
0\le a_i\le C/q_i,\qquad
0\le w_{kj}\le\lambda_k\le C/\theta_k,\\
\ell_k(u_r)\le C/q_i+(C/\theta_k)\sum_j\Delta_{ir,j}
\quad(i\in I_k,\ r\in[N],\ k=0,1).
\end{gathered}
\]
若 $\ell$ 非常数，凸性及单调性给出
$\ell(t)\to\infty$ 当 $t\to\infty$，故两类损失从两侧界住 $u$。
若 $\ell$ 为常数，取 $u=0$。
因此可将原问题限制在一个含有最小值点的紧集中。
若最小值为零，$a=\lambda=w=0$ 将给出
\[
0<2\ell(0)\le\ell(-u_r)+\ell(u_r)\le0,
\]
矛盾。由此得到式~\eqref{eq:appendix-normalized} 及原问题最小值的存在性；
二元等价性进一步给出式~\eqref{eq:appendix-fixed} 最小值的存在性。

\textup{(ii)} 集合 $\mathcal D$ 为锥，且对 $\tau\ge0$，
有 $H(\tau\Pi)=\tau H(\Pi)\ge0$。
对可行三元组，令 $b=\alpha+\mu^\top z$。
若 $b>0$，由归一化及 \textup{(i)} 得 $H(\Pi)/b\le F(z)$，故
\begin{equation}\label{eq:scaling}
2\sqrt{H(\Pi)}-b
=\frac{H(\Pi)}b-
\left(\sqrt{\frac{H(\Pi)}b}-\sqrt b\right)^2\le F(z).
\end{equation}
若 $b=0$，则 $\alpha=0$、$\Pi=0$，仿射目标为零。
反之，任意归一化最优解满足
\[
H(F(z)\bar\Pi)=F(z)^2,\qquad
F(z)(\bar\alpha+\bar\mu^\top z)=F(z),\qquad
2\sqrt{H(F(z)\bar\Pi)}-F(z)=F(z).
\]
因此式~\eqref{eq:appendix-affine} 由式~\eqref{eq:dual-rescaling} 的三元组取得。

\textup{(iii)} 可行域 $\mathcal D$ 不依赖 $z$；
式~\eqref{eq:appendix-affine} 为仿射函数的最大值，故为凸函数。
对偶可行性给出 $c_t(z)\le F(z)$，在 $z^t$ 处的最优性给出
$c_t(z^t)=F(z^t)$。两值相减即得
式~\eqref{eq:support-contact} 与所述次梯度不等式。
归一化乘子必须按
$\mu^t=F(z^t)\bar\mu^t$ 缩放。

对非精确求解，设 $(\Pi,\alpha,\mu)\in\mathcal D$、
$0\le h\le H(\Pi)$，且 $c(z)=2\sqrt h-\alpha-\mu^\top z$。
若 $p_t$ 是 $z^t$ 处的原问题可行目标值，且 $p_t-c(z^t)\le\delta$，则
\[
\begin{gathered}
c(z)\le F(z),\qquad
0\le F(z^t)-c(z^t)\le p_t-c(z^t)\le\delta,\\
F(z)\ge F(z^t)-\mu^\top(z-z^t)-\delta
\quad(z\in[0,1]^D).
\end{gathered}
\]
\end{proof}


\section{收敛性的证明}\label{app:convergence}
设 $x_i\in\mathbb R^D$，非空集合 $I_0,I_1$ 划分 $[N]$，
$q_i,\theta_k>0$ 且 $\sum_{i\in I_k}q_i=1$。
损失 $\ell:\mathbb R\to[0,\infty)$ 有限、凸且非减，满足 $\ell(0)>0$。设 $d\in\{0,\ldots,D\}$。
记

\[
F(z)=\max_{(\Pi,\alpha,\mu)\in\mathcal D}
\{2\sqrt{H(\Pi)}-\alpha-\mu^\top z\},\qquad
\mathcal Z_d=\{z\in\{0,1\}^D:\mathbf1^\top z\le d\},\qquad
\mathrm{OPT}=\min_{z\in\mathcal Z_d}F(z),
\]

其中 $H$、$\mathcal D$ 的定义见附录~\ref{app:dual}。
记 $\Delta_{ir,j}=|x_{ij}-x_{rj}|$，令

\[
C=2\ell(0),\qquad
A_{kj}=\sum_{i\in I_k}q_i\max_{r\in[N]}\Delta_{ir,j},\qquad
M_j=C\max_{k=0,1}\frac{A_{kj}}{\theta_k}.
\]

令 $J_m$ 为 $M_j$ 最大的 $m$ 个坐标的指标集，同值按指标排序。
对 $M_{(1)}\ge\cdots\ge M_{(D)}$，定义

\[
\rho_m=\sum_{r=m+1}^{\min(D,m+d)}M_{(r)},\qquad
K_m=\sum_{r=0}^{\min(d,m)}\binom mr,\qquad 0\le m\le D.
\]

固定 $\delta,\beta\ge0$ 和 $\varepsilon>0$。
在每个已评估点 $z^s=\mathbf1_{S_s}$，取原始可行值 $p_s$，
并恢复 $f_s=g_s\circ\operatorname{Proj}_{S_s}$，故 $R(f_s)\le p_s$。
对应的仿射下界为

\[
\begin{gathered}
c_s(z)=2\sqrt{h_s}-\alpha_s-(\mu^s)^\top z,\qquad
(\Pi^s,\alpha_s,\mu^s)\in\mathcal D,\qquad 0\le h_s\le H(\Pi^s),\\
0\le\mu^s\le M,\qquad 0\le p_s-c_s(z^s)\le\delta.
\end{gathered}
\]

保留全部割，定义

\[
\begin{gathered}
m_t(z)=\max\{0,c_1(z),\ldots,c_t(z)\},\qquad
U_t=\min_{s\le t}p_s,\qquad f_t=f_{s_t},\quad
s_t\in\operatorname*{arg\,min}_{s\le t}p_s,\\
\underline L_t\le\min_{z\in\mathcal Z_d}m_t(z),\qquad
L_t=\max\{0,\underline L_1,\ldots,\underline L_t\},\qquad G_t=U_t-L_t.
\end{gathered}
\]

主问题返回 $z^{t+1}\in\mathcal Z_d$ 和 $\eta^{t+1}$，满足

\[
\eta^{t+1}\ge m_t(z^{t+1}),\qquad
0\le\eta^{t+1}-L_t\le\beta.
\]

\begin{proof}[命题~\ref{prop:convergence} 的证明]
\textup{(i)} 按附录~\ref{app:dual} 将归一化最优解乘以 $F(z)$。
由 $\bar\alpha+\bar\mu^\top z=1$ 及 $F(z)\le C$，有

\[
0\le\alpha=F(z)\bar\alpha\le C,\qquad
d_{kj}(\Pi)=\sum_{i\in I_k,r}\Pi_{ir}\Delta_{ir,j}
\le\alpha\sum_{i\in I_k}q_i\max_r\Delta_{ir,j}=\alpha A_{kj}.
\]

令 $b_j=\max_k d_{kj}(\Pi)/\theta_k$、
$\mu_j'=\min\{\mu_j,b_j\}$。当 $\mu_j>b_j$ 时，

\[
(d_{kj}(\Pi)-\theta_k\mu_j')_+
=(d_{kj}(\Pi)-\theta_k\mu_j)_+=0\quad(k=0,1).
\]

故 $(\Pi,\alpha,\mu')\in\mathcal D$，且

\[
\begin{gathered}
F(z)\ge2\sqrt{H(\Pi)}-\alpha-(\mu')^\top z
\ge2\sqrt{H(\Pi)}-\alpha-\mu^\top z=F(z),\\
0\le\mu_j'\le b_j\le\alpha\max_k A_{kj}/\theta_k\le M_j.
\end{gathered}
\]

因此加入这些界后仍保留最优解及最优值。

\textup{(ii)} 由原始可行性、对偶可行性与 $h_s\le H(\Pi^s)$，得

\begin{equation}\label{eq:inexact-certificate}
0\le m_t\le F,\qquad
L_t\le\min_{z\in\mathcal Z_d}m_t(z)\le\mathrm{OPT}
\le R(f_t)\le U_t.
\end{equation}

由于保留全部割，$m_t$、$L_t$ 非减，而 $U_t$、$G_t$ 非增。
对 $s<r\le t+1$，主问题选点 $z^r=\mathbf1_{S_r}$ 满足

\begin{equation}\label{eq:quantitative-separation}
\begin{aligned}
G_{r-1}
&\le p_s-c_s(z^r)+\eta^r-L_{r-1}\\
&=p_s-c_s(z^s)+(\mu^s)^\top(z^r-z^s)+\eta^r-L_{r-1}\\
&\le\delta+\beta+\sum_{j\in S_r\setminus S_s}M_j.
\end{aligned}
\end{equation}

固定 $m$，假设 $G_t>\delta+\beta+\rho_m$。
由单调性及 \eqref{eq:quantitative-separation}，得

\[
\sum_{j\in S_r\setminus S_s}M_j>\rho_m\qquad(s<r\le t+1).
\]

若 $S_r\cap J_m=S_s\cap J_m$，则
$S_r\setminus S_s\subseteq[D]\setminus J_m$ 且 $|S_r\setminus S_s|\le d$，故

\[
\sum_{j\in S_r\setminus S_s}M_j
\le\sum_{q=m+1}^{\min(D,m+d)}M_{(q)}=\rho_m,
\]

矛盾。因此 $S_1\cap J_m,\ldots,S_{t+1}\cap J_m$ 两两不同，
而这样的子集至多有 $K_m$ 个，故 $t+1\le K_m$。
当 $t\ge K_m$ 时矛盾，从而

\[
0\le R(f_t)-\mathrm{OPT}\le G_t\le\delta+\beta+\rho_m
\qquad(t\ge K_m).
\]

对 $\varepsilon\ge\delta+\beta$，令
$m_\varepsilon=\min\{m\in\{0,\ldots,D\}:\rho_m\le\varepsilon-\delta-\beta\}$。
即得 $G_{K_{m_\varepsilon}}\le\varepsilon$ 及 $T_\varepsilon\le K_{m_\varepsilon}$，
除非算法已经提前停止。计数包含最后一个主问题候选点，无须再进行一次内层求解。特别地，

\[
\rho_D=0,\quad K_D=\sum_{r=0}^d\binom Dr,\qquad
d=0\ \text{or}\ M=0\quad\Longrightarrow\quad G_1\le\delta+\beta.
\]

\end{proof}
\end{document}
